\scanpage{39}
\diagram{\begin{tikzcd}X\arrow[d,"f"]\\Y=\operatorname{Spec}(\unsure{O|\mathcal O})\end{tikzcd}}
\[
F=\widetilde M
\]
$f$ \inserted{plat et} propre, $Y$ schéma noethérien local, $F$ cohérent sur $Y$.
\[
F'=f^*(F),\qquad F'_1\subset F'
\]
\inserted{ss faisceau cohérent.}
\deleted{\unsure{on}} Je dis qu'il y a un plus \unsure{grand|petit} ss-faisceau cohérent $F_1$ de $\unsure{F'|F}$ t.q. $f^*(F_1)\supset F'_1$.

Si $f^*(F_1)$ et $f^*(F_2)\supset F'_1$ alors
\[
f^*(F_1\cap F_2)=f^*(F_1)\cap f^*(F_2)\supset F'_1
\]
\margin{platitude de $f$}
$F_i\subset F$ tq $f^*(F_i)\supset F'_1$ \deleted{\unsure{forment}} \unsure{une famille filtrante décroissante}. Il reste à prouver que $F_1=\bigcap F_i$, alors $f^*(F_1)\supset F'_1$ i.e. l'image de $F'_1$ dans $F'/f^*(F_1)=f^*(F/F_1)$ est nulle. Remplaçant $F$ par $F/F_1$, on est ramené à prouver ceci : \uline{si} $\bigcap F_i=0$, \uline{alors} \deleted{$\bigcap\unsure{f^*(F_i)}=0$} \inserted{tout faisceau cohérent $F'_1$ contenu dans $\bigcap f^*(F_i)$ est nul} [$f$ propre \deleted{\unsure{complet|surjectif}}, les $F_i$ filtrants décr. dans $F$ cohérent, $Y$ local et complet].
\margin{[\unsure{on sait même que} $\bigcap f^*(F_i)$ est nul ?]}
\deleted{Comme $F'_1=\bigcap\unsure{f^*(F_i)}$ \unsure{de}}
\deleted{Soit $s$ une section de}

$F'_1$ est fermé, il suffit de prouver qu'il ne rencontre pas $f^{-1}(y)$ ($y$ le pt fermé de $Y$) donc de prouver que si $s$ est une section de $F'_1$ sur un ouvert $U$, elle est nulle en $\forall\ x\in f^{-1}(y)$. Or comme $\bigcap M_i=0$, les $M_i$ tendent vers $0$ dans $M$ (?)

\margin{N.B. Faut définir \unsure{pour un} $\unsure{O|\mathcal O}$-module $M$, les $F_i$ par des $M_i\subset M$, pour \uline{\unsure{définition}} \unsure{en vrai sens}, $\unsure{\bigcap_iF_i}=\widetilde{\bigcap M_i}$ (ce n'est pas le faisceau $\bigcap$, mais \unsure{l'inf dans les faisceaux qu. coh.}).}
\margin{Énoncé : \unsure{ss les|sous les} conditions envisagées ($f$ propre, les $F_i$ décroissants dans $F$ cohérent, $Y$ local et complet), $f^*(\bigcap F_i)=\bigcap f^*(F_i)$ ?}
\margin{\uline{\unsure{Arranger les notations}}}
% archive foliation: 38
